\documentclass[11pt,reqno]{amsart}

\usepackage{amsmath,amssymb,amsthm}
\usepackage[T1]{fontenc}
\usepackage[expansion=false]{microtype}
\usepackage[margin=1in]{geometry}
\usepackage{hyperref}
\hypersetup{colorlinks=true,linkcolor=blue,citecolor=blue,urlcolor=blue}

\theoremstyle{plain}
\newtheorem{theorem}{Theorem}[section]
\newtheorem{proposition}[theorem]{Proposition}
\newtheorem{lemma}[theorem]{Lemma}
\newtheorem{corollary}[theorem]{Corollary}

\theoremstyle{definition}
\newtheorem{definition}[theorem]{Definition}
\newtheorem{observation}[theorem]{Observation}
\newtheorem{remark}[theorem]{Remark}

\newcommand{\Q}{\mathbb{Q}}
\newcommand{\Z}{\mathbb{Z}}
\newcommand{\F}{\mathbb{F}}
\newcommand{\R}{\mathbb{R}}
\newcommand{\PP}{\mathbb{P}}
\newcommand{\Aut}{\operatorname{Aut}}
\newcommand{\Bir}{\operatorname{Bir}}
\DeclareMathOperator{\Jac}{Jac}
\DeclareMathOperator{\rank}{rank}

\sloppy

\begin{document}

\title[No perfect cuboid in Case~B at \texorpdfstring{$p=1$}{p=1}]{No perfect cuboid in Case~B at \texorpdfstring{$p=1$}{p=1}, and the rank obstruction on the associated genus-five curve}

\author{Lightman Chang}
\address{Independent Researcher}
\email{lightman.chang@gmail.com}

\subjclass[2020]{Primary 11D09; Secondary 11G05, 11G30, 14H45}
\keywords{perfect cuboid, Euler brick, Pell equation, Lucas sequence, genus-five curve, Mordell--Weil rank, Chabauty--Coleman}

\date{\today}

\begin{abstract}
A perfect cuboid is a rectangular box whose three edges, three face diagonals,
and space diagonal are all integers; whether one exists is a question of Euler
(1769) that remains open. Inside the standard two-adic case analysis of
primitive perfect cuboids, the stratum ``Case~B at $p=1$'' is the
one-parameter family with edges $(4q,\,q^2-4,\,2(q^2-1))$, $q$ a positive
integer; the space diagonal is integral exactly when $g^2=5q^4+20$. We prove
that this stratum contains no perfect cuboid: setting $Y=q^2$ turns the space
condition into the Pell equation $g^2-5Y^2=20$, whose nonnegative solutions
$Y_n$ are the odd-indexed Lucas numbers $L_{2n-1}$, and a theorem of Cohn
(1964) on squares in the Lucas sequence forces $Y\in\{1,4\}$, i.e.\ $q\in\{1,2\}$,
both degenerate. The argument is unconditional and entirely verified in
PARI/GP. We then analyse the associated genus-five curve
$C:\{e^2=5q^4-16q^2+20,\ g^2=5q^4+20\}$, whose finitely many rational points
encode all \emph{rational} parameters and which Peschmann (2026) singled out as
the target of a prospective genus-five Chabauty--Coleman obstruction. We
determine the $\Q$-isogeny decomposition of $\Jac(C)$ explicitly and show that
all five elliptic factors have Mordell--Weil rank one, so
$\rank\Jac(C)(\Q)=5=\operatorname{genus}(C)$. Chabauty--Coleman therefore does
\emph{not} apply to $C$; the factors carrying the relevant differentials are the
quadratic twists by $\Q(\sqrt5)$ of two rank-zero curves, not those curves
themselves. This isolates the elementary Pell--Lucas argument as the one that
closes the stratum. A closing section records why elementary descent also stalls:
the perfect-cuboid complete intersection $V\subset\PP^6$ is a surface of general
type with finite $\Bir(V)$, an obstruction to any birational height-strict descent.
\end{abstract}

\maketitle

\section{Introduction}\label{sec:intro}

A \emph{perfect cuboid} is a rectangular parallelepiped with positive integer
edges $a,b,c$ whose three face diagonals $\sqrt{a^2+b^2}$, $\sqrt{b^2+c^2}$,
$\sqrt{a^2+c^2}$ and whose space diagonal $\sqrt{a^2+b^2+c^2}$ are all integers.
The Perfect Cuboid Problem (PCP), which asks whether a perfect cuboid exists,
was raised by Euler in 1769 and remains open after more than two and a half
centuries; no perfect cuboid is known, and exhaustive searches have excluded
any perfect cuboid with smallest edge below $5\times10^{11}$ or odd edge below
$2.5\times10^{13}$~\cite{matson}. Since no
parametrization of all Euler bricks is available, a standard strategy is to
stratify the primitive solutions by their two-adic structure and to dispatch
explicitly parametrized strata one at a time.

Two recent preprints frame the elliptic-curve geometry that surrounds these
strata. Yoshida~\cite{yoshida} studies \emph{face cuboids} through the family
$E_{1,s}:y^2=x(x-(2s)^2)(x+(s^2-1)^2)$ and uses non-torsion points to
\emph{produce} infinitely many of them. Peschmann~\cite{peschmann} reduces the
full perfect-cuboid problem to a one-parameter family of genus-three curves
$C_A:w^2=\lambda^8+A\lambda^4+1$ and develops $2$-descent and Kummer-character
obstructions. In his closing discussion, Peschmann introduces a genus-five
double cover $C_{T_4}:z^2=f(\,\cdot+T_4)/f(\,\cdot)$ of an elliptic quotient
and observes that \emph{if} $\rank\Jac(C_{T_4})(\Q)<5$ then Chabauty--Coleman
would render the relevant finiteness effective, while cautioning that ``for
larger parameters, rank sums exceeding $5$ cannot be excluded''
(\cite[\S8]{peschmann}). The present paper concerns a concrete genus-five curve
of exactly this shape and shows that, for it, the rank does reach the genus, so
the prospective obstruction is unavailable; the stratum is instead closed by an
elementary argument.

We work in the stratum we call \emph{Case~B at $p=1$}. To fix conventions: in
the standard $2$-adic case analysis of primitive perfect cuboids one stratifies
candidates by the $2$-adic valuations of the edges and of the relevant
face-diagonal data; the present stratum is the one obtained by taking the
distinguished even edge $a$ with $v_2(a)=2$ (i.e.\ $a=4q$ with $q$ odd in the
fully primitive normalization) and specializing the residual integer parameter,
here written $p$, to its base value $p=1$. With these choices the candidate is
the one-parameter family
\begin{equation}\label{eq:caseB}
(a,b,c)=\bigl(4q,\;q^2-4,\;2(q^2-1)\bigr),\qquad q\in\Z_{>0},
\end{equation}
for which the three face conditions hold identically and the space-diagonal
condition is $g^2=a^2+b^2+c^2=5q^4+20$. We adopt this parametrization as arising
from the standard $2$-adic case analysis. The full taxonomy of the remaining
strata is treated in the companion analysis; it is not standard in the
published literature, and we make no claim that it is. Our main result is that
this stratum contains no perfect cuboid (Theorem~\ref{thm:main}); the proof
reduces the space condition to a Pell equation whose square values are governed
by Cohn's theorem~\cite{cohn,cohnlucas} on squares in the Lucas sequence, and it
is unconditional. We then turn to the genus-five curve
$C:\{e^2=5q^4-16q^2+20,\ g^2=5q^4+20\}$ that
records \emph{all rational} parameters, give the explicit $\Q$-isogeny
decomposition of its Jacobian, and prove that every elliptic factor has rank
one, so $\rank\Jac(C)(\Q)=5=\operatorname{genus}(C)$
(Proposition~\ref{prop:rank}). This identifies the precise obstruction to the
curve-theoretic route on $C$ and corrects a natural but mistaken expectation:
the two factors carrying the candidate Chabauty differentials are the
$\Q(\sqrt5)$-quadratic twists of the rank-zero curves \texttt{120a2} and
\texttt{80a1}, and the twists have rank one.

The paper is organized as follows. Section~\ref{sec:closure} sets up Case~B at
$p=1$, derives the Pell equation, identifies its orbit with the odd-indexed
Lucas numbers, and applies Cohn's theorem to prove the main theorem.
Section~\ref{sec:curve} introduces the genus-five curve, computes the
$\Q$-isogeny decomposition of its Jacobian and the rank of each factor, and
records the resulting rank obstruction together with the unique genus-two
quotient. Section~\ref{sec:scope} delimits the scope of the result---a single
stratum, not the full problem---and positions it relative to Peschmann and
Yoshida. Section~\ref{sec:obstructions} then records structural obstructions
that situate the elementary route within the wider geometry: the perfect-cuboid
complete intersection $V\subset\PP^6$ is a minimal surface of general type, so
its birational automorphism group is finite and no birational self-map of $V$
can effect a height-strict Markov--Vieta or Fermat descent
(Proposition~\ref{prop:descent}); and the standard per-fiber elliptic curve has
purely multiplicative reduction, which obstructs an elementary discriminant-free
height lower bound (Observation~\ref{obs:noabs}). These are honest negatives,
delimiting why elementary descent stalls without bearing on whether a perfect
cuboid exists. All arithmetic claims are verified in PARI/GP, with scripts and
captured output accompanying the paper.

\section{Case~B at \texorpdfstring{$p=1$}{p=1} and its closure}\label{sec:closure}

\subsection{The stratum and its reduction to a Pell equation}

\begin{definition}\label{def:caseB}
For a positive integer $q$, the \emph{Case~B at $p=1$ candidate} is the triple
\[
\mathrm{B}(q):=\bigl(a,b,c\bigr)=\bigl(4q,\;q^2-4,\;2(q^2-1)\bigr).
\]
It is a \emph{perfect cuboid} if all three edges are positive and
$a^2+b^2$, $b^2+c^2$, $a^2+c^2$, $a^2+b^2+c^2$ are all perfect squares.
\end{definition}

\begin{lemma}\label{lem:faces}
For every $q$ the candidate $\mathrm{B}(q)$ satisfies the polynomial identities
\[
a^2+b^2=(q^2+4)^2,\qquad a^2+c^2=\bigl(2(q^2+1)\bigr)^2,
\]
\[
b^2+c^2=5q^4-16q^2+20,\qquad a^2+b^2+c^2=5q^4+20 .
\]
Consequently the two face diagonals $\sqrt{a^2+b^2}=q^2+4$ and
$\sqrt{a^2+c^2}=2(q^2+1)$ are automatically integral, and $\mathrm{B}(q)$ is a
perfect cuboid if and only if $b^2+c^2$ and $a^2+b^2+c^2$ are simultaneously
perfect squares and $b=q^2-4>0$.
\end{lemma}

\begin{proof}
The four displayed equalities are identities in $\Z[q]$, verified by direct
expansion in script~\texttt{08\_caseB\_parametrization.gp}: with $a=4q$,
$b=q^2-4$ and $c=2(q^2-1)$, one has $a^2+b^2-(q^2+4)^2=0$ and
$a^2+c^2-4(q^2+1)^2=0$ identically, while $b^2+c^2=5q^4-16q^2+20$ and
$a^2+b^2+c^2=5q^4+20$. The first two identities make the corresponding face
diagonals integral for all integers $q$; the remaining two conditions, together
with positivity of the edges, are then exactly what is needed for a perfect
cuboid. Positivity of $a=4q$ and $c=2(q^2-1)$ holds for every $q\ge2$, so the
binding positivity constraint is the requirement $b=q^2-4>0$, that is, $q\ge3$.
\end{proof}

\begin{lemma}\label{lem:pell}
Write $g^2=a^2+b^2+c^2$ and set $Y=q^2$. Then the space-diagonal condition
$g^2=5q^4+20$ is equivalent to the Pell-type equation
\[
g^2-5\,Y^2=20,\qquad Y=q^2 .
\]
Its nonnegative integer solutions $(g,Y)$ form a single orbit under the totally
positive unit $\tfrac{3+\sqrt5}{2}$ of $\Z\!\left[\tfrac{1+\sqrt5}{2}\right]$,
with $Y$-coordinates $Y_1,Y_2,\dots$ generated by
\[
Y_1=1,\quad Y_2=4,\quad Y_{n+1}=3Y_n-Y_{n-1}.
\]
\end{lemma}

\begin{proof}
Substituting $Y=q^2$ into $g^2=5q^4+20$ gives $g^2-5Y^2=20$ directly. The
fundamental solution is $(g,Y)=(5,1)$, since $25-5=20$. Multiplication of
$g+Y\sqrt5$ by the totally positive unit $\tfrac{3+\sqrt5}{2}$ of norm $1$ sends
solutions to solutions and produces the recurrence $g_{n+1}=3g_n-g_{n-1}$,
$Y_{n+1}=3Y_n-Y_{n-1}$ with $(g_1,Y_1)=(5,1)$ and $(g_2,Y_2)=(10,4)$; the first
ten terms of $Y_n$ are
$1,4,11,29,76,199,521,1364,3571,9349$
(script~\texttt{07\_known\_points\_and\_cohn.gp}).

We now show that every nonnegative solution lies in this single orbit, by
reducing $g^2-5Y^2=20$ to a classical negative-Pell equation. Since $5\mid 20$
and $5\mid 5Y^2$, the relation $g^2=5Y^2+20$ forces $5\mid g^2$, hence $5\mid g$.
(Note $5\nmid Y$: if $5\mid Y$ then $25\mid 5Y^2$ and $g^2\equiv 20\pmod{25}$,
impossible since $20$ is not a square modulo $25$.) Writing $g=5h$ gives
$25h^2-5Y^2=20$, i.e.
\[
Y^2-5h^2=-4 .
\]
The positive solutions of $X^2-5Z^2=-4$ are classically known to be exactly
$(X,Z)=(L_{2n-1},F_{2n-1})$, $n\ge1$, where $L_k$ and $F_k$ are the Lucas and
Fibonacci numbers; these form a single solution class, the negative-Pell
equation associated with the Lucas/Fibonacci identity $L_m^2-5F_m^2=4(-1)^m$
\cite[Ch.~7]{niven}. Hence $Y\in\{L_{2n-1}:n\ge1\}$, and correspondingly
$h=F_{2n-1}$, $g=5h=5F_{2n-1}$, recovering precisely the single orbit above.
(For the first few terms, $(Y,h)=(1,1),(4,2),(11,5)$ satisfy
$1-5=-4$, $16-20=-4$, $121-125=-4$.) Thus every nonnegative solution of
$g^2-5Y^2=20$ has $Y=L_{2n-1}$ for some $n\ge1$.
\end{proof}

\begin{lemma}\label{lem:lucas}
The sequence $Y_n$ of Lemma~\ref{lem:pell} satisfies $Y_n=L_{2n-1}$, where
$L_k$ is the Lucas sequence $L_1=1$, $L_2=3$, $L_{k+1}=L_k+L_{k-1}$.
\end{lemma}

\begin{proof}
The Lucas sequence satisfies $L_{2n+1}=3L_{2n-1}-L_{2n-3}$ (a consequence of
$L_{k+2}=L_{k+1}+L_k$ applied twice), with $L_1=1$ and $L_3=4$. Thus
$L_{2n-1}$ obeys the same second-order recurrence and the same two initial
values as $Y_n$, so the two sequences coincide. The equality $Y_n=L_{2n-1}$ is
confirmed for $1\le n\le 12$ in script~\texttt{07}.
\end{proof}

\subsection{The main theorem}

\begin{theorem}\label{thm:main}
No Case~B at $p=1$ candidate is a perfect cuboid. That is, for every positive
integer $q$ the triple $\mathrm{B}(q)=\bigl(4q,\,q^2-4,\,2(q^2-1)\bigr)$ fails
to be a perfect cuboid: either some edge is nonpositive, or $a^2+b^2+c^2$ is not
a perfect square.
\end{theorem}

\begin{proof}
Suppose $\mathrm{B}(q)$ is a perfect cuboid. By Lemma~\ref{lem:faces} this
requires $a^2+b^2+c^2=5q^4+20$ to be a perfect square, say $g^2$, and the edge
$b=q^2-4$ to be positive, i.e.\ $q\ge3$. By Lemma~\ref{lem:pell}, writing
$Y=q^2$, the value $Y$ must be a member $Y_n$ of the Pell orbit; in particular
$q^2=Y_n$ is a perfect square. By Lemma~\ref{lem:lucas}, $Y_n=L_{2n-1}$, so a
solution requires the Lucas number $L_{2n-1}$ to be a perfect square.

By Cohn's theorem on squares in the Lucas sequence~\cite{cohn,cohnlucas} (the
Lucas-square statement is the one established in
\cite{cohnlucas}), the only perfect
squares among the Lucas numbers $L_k$ ($k\ge1$) are $L_1=1$ and $L_3=4$. Hence
$Y_n=L_{2n-1}\in\{1,4\}$, giving $q^2\in\{1,4\}$, i.e.\ $q\in\{1,2\}$. For $q=1$
the edge $b=q^2-4=-3$ is negative (and $c=2(q^2-1)=0$ vanishes); for $q=2$ the
edge $b=q^2-4=0$ vanishes. In both cases the candidate is degenerate, so no
$q\ge3$ yields a perfect cuboid. This contradiction proves the theorem; the
arithmetic is verified in script~\texttt{08} (face and space identities) and
script~\texttt{07} (Pell orbit, Lucas identification, and the square test
returning $\{1,4\}$).
\end{proof}

\begin{remark}
The closure is unconditional and elementary: it rests only on the explicit
parametrization~\eqref{eq:caseB}, the structure of solutions to a Pell equation,
and Cohn's theorem~\cite{cohn,cohnlucas}, all of which are unconditional. It does
\emph{not} invoke the Mordell conjecture, the method of Chabauty, or any
hypothesis on $L$-functions. The integrality of the parameter $q$---which is
what makes the Pell--Lucas argument available---is built into the
Pythagorean-style parametrization~\eqref{eq:caseB}: a Case~B at $p=1$ cuboid has
integer edges $a=4q$, so $q$ is an integer.
\end{remark}

\section{The genus-five curve and the rank obstruction}\label{sec:curve}

The stratum of Section~\ref{sec:closure} is parametrized by integers $q$. It is
natural to ask the stronger question of which \emph{rational} $q$ make both
$5q^4-16q^2+20$ and $5q^4+20$ squares; these are the rational points of the
joint curve
\begin{equation}\label{eq:C}
C:\qquad e^2=5q^4-16q^2+20,\qquad g^2=5q^4+20 ,
\end{equation}
the fibre product over $\mathbb{P}^1_q$ of the two genus-one double covers
$e^2=f_1(q)$ and $g^2=f_2(q)$, with $f_1=5q^4-16q^2+20$ and $f_2=5q^4+20$. This
is the curve singled out, in abstract form, by Peschmann~\cite[\S8]{peschmann}
as the locus on which an effective genus-five Chabauty--Coleman argument might
operate. We show in this section that the rank obstruction he flagged is
present for $C$.

\begin{proposition}\label{prop:genus}
The polynomials $f_1,f_2$ are coprime and separable, neither has a rational
root, and the smooth projective model of $C$ has genus $5$. Its rational points
are affine: the four geometric points at infinity are defined over
$\Q(\sqrt5)$, hence none is rational.
\end{proposition}

\begin{proof}
One computes $\operatorname{res}(f_1,f_2)=2^{20}5^4\neq0$, so $f_1,f_2$ are
coprime, and $\operatorname{disc}(f_1)=33177600$,
$\operatorname{disc}(f_2)=256000000$ are nonzero, so both are separable; neither
has a rational root (their roots are not real, and a rational root would divide
$20$ but none does). The covering $C\to\mathbb{P}^1_q$ has degree $4$. Over each
of the four roots of $f_1$ the $e$-cover ramifies and the $g$-cover does not,
contributing $2$ to the ramification divisor; symmetrically for the four roots
of $f_2$; over $q=\infty$ the cover is unramified. Riemann--Hurwitz gives
$2g(C)-2=4(2\cdot0-2)+16=8$, so $g(C)=5$. At infinity, in the chart
$q=1/s$, $E=es^2$, $G=gs^2$, the equations become $E^2=5-16s^2+20s^4$,
$G^2=5+20s^4$, with $E^2=G^2=5$ at $s=0$; since $5$ is not a rational square the
four points at infinity are defined over $\Q(\sqrt5)$. These computations are in
script~\texttt{01\_genus\_and\_decomposition.gp}.
\end{proof}

\begin{proposition}[$\Q$-isogeny decomposition of $\Jac(C)$]\label{prop:decomp}
The curve $C$ carries the group $(\Z/2\Z)^3$ of involutions generated by
$\sigma_e:e\mapsto-e$, $\sigma_g:g\mapsto-g$, $\sigma_q:q\mapsto-q$, and there
is a $\Q$-isogeny
\[
\Jac(C)\;\sim_{\Q}\;E_1\times E_2\times E_3\times X_+^{(5)}\times X_-^{(5)},
\]
where $E_1,E_2,E_3$ are the elliptic curves of Cremona labels \texttt{480f1},
\texttt{800a1}, \texttt{1200a2} (the quotients by $\sigma_g$, $\sigma_e$ and a
factor of the quotient by $\sigma_q$), and $X_+^{(5)},X_-^{(5)}$ are the
quadratic twists by $\Q(\sqrt5)$ of $X_+=\texttt{120a2}$ and $X_-=\texttt{80a1}$,
with Cremona labels \texttt{600a2} and \texttt{400a2}.
\end{proposition}

\begin{proof}
The two ``axis'' quotients $C/\langle\sigma_g\rangle:e^2=f_1$ and
$C/\langle\sigma_e\rangle:g^2=f_2$ have, after minimalization, the models
$y^2=x^3-39312x+2889216$ ($j=48228544/2025$, conductor $480$) and
$y^2=x^3-32400x$ ($j=1728$, CM by $\Z[i]$, conductor $800$); PARI's
\texttt{ellfromeqn} applied to $e^2-f_1$ and $g^2-f_2$ reproduces these
$j$-invariants and conductors, identifying them as \texttt{480f1} and
\texttt{800a1} (script~\texttt{01}). The third factor $E_3=\texttt{1200a2}$,
model $y^2=x^3-x^2-108x-288$, arises from the genus-two quotient by $\sigma_q$.

The remaining two factors carry the $(-1,-1,-1)$-eigenspace of the holomorphic
differentials, spanned by $\omega_1=dq/(eg)$ and $\omega_3=q^2\,dq/(eg)$. We
identify them by comparing the Frobenius data of $C$ with the candidate
factors. Counting points of the smooth model of $C$ over $\F_p$ and forming
$a_p(C)=p+1-\#C(\F_p)$, one finds
\[
a_p(C)=a_p(E_1)+a_p(E_2)+a_p(E_3)+a_p\!\bigl(X_+^{(5)}\bigr)+a_p\!\bigl(X_-^{(5)}\bigr)
\]
for every good prime $p$ in $[7,150]$ ($32$ primes, no exceptions), where
$X_+^{(5)},X_-^{(5)}$ are the twists by $5$ of \texttt{120a2}, \texttt{80a1};
PARI's \texttt{ellidentify} returns the labels \texttt{600a2}, \texttt{400a2}
for these twists (script~\texttt{06\_true\_decomposition\_rank.gp}). Matching
the untwisted curves \texttt{120a2}, \texttt{80a1} instead fails at every prime
with $\bigl(\tfrac5p\bigr)=-1$; the discrepancy is exactly the sign of
$\bigl(\tfrac5p\bigr)$, the signature of a $\Q(\sqrt5)$-twist (scripts
\texttt{01}, \texttt{02\_genus2\_quotient.gp}). The agreement of all five linear
Frobenius coefficients over this range of $32$ primes, taken together with the
explicit quotient constructions of scripts~\texttt{04} and~\texttt{05} and the
differential-character decomposition recorded in
Proposition~\ref{prop:g2quot}, identifies the isogeny-class decomposition of
$\Jac(C)$ displayed above; the $a_p$ agreement furnishes strong corroborating
evidence, and it is the explicit constructions that pin down the factors. (We
do not invoke a Faltings--Serre argument here, as that would require a specific
computable bound on the primes to be checked.)
\end{proof}

\begin{proposition}[The rank obstruction]\label{prop:rank}
Each of the five factors $E_1=\texttt{480f1}$, $E_2=\texttt{800a1}$,
$E_3=\texttt{1200a2}$, $X_+^{(5)}=\texttt{600a2}$, $X_-^{(5)}=\texttt{400a2}$
has Mordell--Weil rank one over $\Q$. Consequently
\[
\rank\Jac(C)(\Q)=1+1+1+1+1=5=\operatorname{genus}(C),
\]
so the Chabauty hypothesis $\rank<\operatorname{genus}$ fails and the method of
Chabauty--Coleman does not apply to $C$.
\end{proposition}

\begin{proof}
For each of the five curves PARI's \texttt{ellrank}---an unconditional
$2$-descent computation---returns the tight interval $[1,1]$ together with an
explicit point of infinite order, namely $(-48,2160)$ on $E_1$, $(-144,1296)$
on $E_2$, $(-4,8)$ on $E_3$, $(-16,72)$ on $X_+^{(5)}$, and $(1,24)$ on
$X_-^{(5)}$; each listed point has nonzero canonical height and infinite order
(scripts \texttt{05\_genuine\_quotient\_rank.gp}, \texttt{06}). Independently,
each curve has analytic rank one with root number $-1$, so Kolyvagin's
theorem~\cite{kolyvagin} also forces algebraic rank one. The five ranks sum to
$5$, equal to the genus by Proposition~\ref{prop:genus}, and Chabauty's
inequality requires the strict inequality $\rank\Jac(C)(\Q)<\operatorname{genus}(C)$.
\end{proof}

\begin{remark}\label{rem:wrongfactors}
Proposition~\ref{prop:rank} corrects a natural misidentification. The two
differentials $\omega_1,\omega_3$ are pulled back from the
$(-1,-1,-1)$-eigenspace, and one might expect this eigenspace to be governed by
the rank-zero curves $X_+=\texttt{120a2}$ and $X_-=\texttt{80a1}$ (both
unconditionally of rank $0$: PARI's \texttt{ellrank} gives $[0,0]$, and
$L(X_+,1)\approx1.2695$, $L(X_-,1)\approx1.0095$ are nonzero, so Kolyvagin
applies). If those were the actual factors, $\rank\Jac(C)(\Q)$ would be $3$ and
Chabauty would apply. But the factors are their $\Q(\sqrt5)$-twists
$\texttt{600a2}$ and $\texttt{400a2}$---a twist forced by the appearance of
$\sqrt5$ at the points at infinity (Proposition~\ref{prop:genus})---and the
twists each have rank one. Rank is not a quasi-isogeny invariant under
quadratic twist, and here the twist raises the rank from $0$ to $1$ in each
factor.
\end{remark}

\begin{proposition}[The unique genus-two quotient]\label{prop:g2quot}
Among the quotients of $C$ by index-two subgroups of $(\Z/2\Z)^3$, exactly one
has genus two, namely $C_q:=C/H$ with $H=\langle\sigma_e\sigma_g,\,
\sigma_e\sigma_q\rangle$, whose invariant differentials are $\omega_1,\omega_3$.
A plane model is
\[
C_q:\qquad w^2=u\,(5u^2-16u+20)(5u^2+20),\qquad u=q^2,\ w=egq,
\]
and $\Jac(C_q)\sim_{\Q}X_+^{(5)}\times X_-^{(5)}=\texttt{600a2}\times\texttt{400a2}$.
In particular $\rank\Jac(C_q)(\Q)=2=\operatorname{genus}(C_q)$, so Chabauty does
not apply to $C_q$ either.
\end{proposition}

\begin{proof}
A holomorphic differential of $C$ with $(\Z/2\Z)^3$-character
$(\chi_e,\chi_g,\chi_q)$ is invariant under an index-two subgroup
$H=\ker\psi$ if and only if its character lies in $\{1,\psi\}$. The five
characters occurring are $(-,+,-),(+,-,-),(-,-,+)$ and the repeated
$(-,-,-)$. Running over the seven nontrivial $\psi$, the only choice giving a
two-dimensional invariant space is $\psi=(-,-,-)$, whose invariant forms are
$\omega_1,\omega_3$; all others give dimension $\le1$. Thus the genus-two
quotient is unique. The functions $u=q^2$ and $w=egq$ are exactly the
$H$-invariants, and $w^2=e^2g^2q^2=f_1(q)f_2(q)\,q^2=u(5u^2-16u+20)(5u^2+20)$,
verified as an identity in script~\texttt{04\_genuine\_Cq\_lift.gp}. Counting
$\F_p$-points of this quintic model gives $a_p(C_q)=a_p(X_+^{(5)})+a_p(X_-^{(5)})$
for all $43$ good primes $p\le200$, and the full Frobenius characteristic
polynomial (both coefficients, via $\#C_q(\F_p)$ and $\#C_q(\F_{p^2})$) factors
as $L_p(X_+^{(5)})L_p(X_-^{(5)})$ for all $12$ good primes in $[7,47]$ (scripts
\texttt{02}, \texttt{03\_Cq\_model\_and\_rank.gp}). By
Proposition~\ref{prop:rank} both factors have rank one, so
$\rank\Jac(C_q)(\Q)=2$.
\end{proof}

\begin{remark}\label{rem:knownpts}
The curve $C$ has $16$ visible rational points
$\{(\pm1,\pm3,\pm5),(\pm2,\pm6,\pm10)\}$ (independent signs), all with
$q\in\{\pm1,\pm2\}$ and hence $b=q^2-4\in\{-3,0\}\le0$, i.e.\ all degenerate; a
height search over $q=n/d$ with $|n|,d\le2000$ finds no others
(script~\texttt{07}). Under the quotient map of
Proposition~\ref{prop:g2quot} these $16$ points collapse to the four points
$(u,w)\in\{(1,\pm15),(4,\pm120)\}$ of $C_q$, which lift back to exactly the $16$
(script~\texttt{04}). By the Mordell conjecture (Faltings) $C(\Q)$ is finite,
but the present rank computation shows this finiteness is not made effective by
Chabauty--Coleman on $C$ or on $C_q$; an effective determination of $C(\Q)$
would require a higher method (for instance quadratic Chabauty, or a Mordell--Weil
sieve combined with covering collections), which we do not carry out here. The
closure of the stratum in Theorem~\ref{thm:main} does not depend on knowing
$C(\Q)$: it uses only the integer points, which the Pell--Lucas argument settles
unconditionally.
\end{remark}

\section{Scope and related work}\label{sec:scope}

\subsection{Scope: a single stratum}\label{sec:scopelimit}

Theorem~\ref{thm:main} closes the stratum ``Case~B at $p=1$,'' a single
one-parameter family inside the two-adic case analysis of primitive perfect
cuboids; it does not resolve the Perfect Cuboid Problem, and must not be read as
doing so. Even within Case~B, the present argument covers only $p=1$: for higher
$p$ the analogous space-diagonal equation $g^2=5q^4-\cdots$ depends on $p$ and
the Pell--Lucas reduction changes, so those sub-cases are not covered here.
Case~A and the other face-III parametrizations, the higher two-adic strata, and
non-primitive sub-cases are likewise outside the scope of this paper. Indeed the
companion question of which \emph{rational} parameters $q$ are admissible---the
rational points of the genus-five curve $C$---is shown in
Section~\ref{sec:curve} to be inaccessible to Chabauty--Coleman, precisely
because $\rank\Jac(C)(\Q)=\operatorname{genus}(C)$. We make no claim about the
full problem, and in particular do not claim that PCP is solved.

\subsection{Relation to Peschmann and Yoshida}\label{sec:related}

Peschmann~\cite{peschmann} reduces the full perfect-cuboid problem to a
genus-three family $C_A:w^2=\lambda^8+A\lambda^4+1$ and, in his closing
discussion, introduces a genus-five double cover $C_{T_4}$ of an elliptic
quotient as a possible site for an effective Chabauty--Coleman obstruction,
explicitly noting that he does not decompose the relevant abelian fourfold and
that ``rank sums exceeding $5$ cannot be excluded'' (\cite[\S8]{peschmann}). Our
curve $C$ of~\eqref{eq:C} is a concrete genus-five curve of this type attached to
Case~B at $p=1$. For it we carry out the decomposition Peschmann left abstract
(Proposition~\ref{prop:decomp}) and find that the rank does reach the genus
(Proposition~\ref{prop:rank}); thus, for this instance, the prospective
Chabauty--Coleman obstruction is unavailable, confirming his stated concern
rather than realizing his hoped-for effectivity. The stratum is nonetheless
closed, by the elementary Pell--Lucas argument of Section~\ref{sec:closure},
which is independent of the curve's rational-point count. We regard this as a
precise, honest answer to the concrete instance underlying his \S8 discussion:
the curve-theoretic route is obstructed here, and finiteness for the cuboid
question comes from elsewhere.

Yoshida~\cite{yoshida} studies face cuboids via the elliptic family $E_{1,s}$
and uses non-torsion points to produce infinitely many of them, an existence
direction opposite to ours; the curves $X_+,X_-$ and their twists that arise
here do not appear in his work. The factorization $5q^4+20=5(q^4+4)$ and the
attendant Gaussian/$\Q(\sqrt5)$ structure are elementary, and the rank-zero
curves \texttt{120a2}, \texttt{80a1} are routine entries in Cremona's
tables~\cite{cremona}; the curve \texttt{80a1} also appears in the Saunderson
reduction of a separate parametric family. To the best of our knowledge, the
explicit closure of Case~B at $p=1$ by the Pell--Lucas route, and the
determination that the associated genus-five Jacobian has rank equal to its
genus---with the $(-1,-1,-1)$-eigenspace governed by the rank-one twists
\texttt{600a2}, \texttt{400a2} rather than the rank-zero curves
\texttt{120a2}, \texttt{80a1}---do not appear in the cited literature. We do
not claim that the underlying rank facts are unknown to specialists; we claim
only that this precise statement and proof do not appear there.

\section{Structural obstructions to elementary descent}\label{sec:obstructions}

The closure of Section~\ref{sec:closure} is elementary, and
Section~\ref{sec:curve} shows that the curve-theoretic route on $C$ is obstructed
by rank. It is natural to ask whether some \emph{other} elementary mechanism---a
Markov--Vieta or Fermat-style infinite descent acting directly on the cuboid
locus---could close more of the problem. This section records two structural
reasons such descents stall. The results are delimitative: they explain why
elementary descent on the ambient does not gain traction, and they make no claim
about deciding whether a perfect cuboid exists.

The natural ambient is the smooth complete intersection of four quadrics
\[
V:\quad a^2+b^2-d^2=b^2+c^2-e^2=a^2+c^2-f^2=a^2+b^2+c^2-g^2=0
\]
in $\PP^6$, whose $\Q$-points with all coordinates nonzero are exactly the perfect
cuboids; forgetting the fourth quadric gives the Euler-brick surface
$V'\subset\PP^5$, a (singular) K3 \cite{vanLuijk}.

\begin{proposition}\label{prop:geomV}
The surface $V$ is a smooth minimal surface of general type with ample canonical
class
\[
K_V=\mathcal{O}_V(1),\qquad K_V^2=16,\quad c_2(V)=80,\quad
\chi(\mathcal{O}_V)=8,\quad p_g=7,\quad q=0 .
\]
In particular $V$ is \emph{not} a K3 surface; only the Euler-brick double cover
$V'$ is K3. Its automorphism group is finite, equal to its birational
automorphism group, and is linear:
\[
\Aut(V)=\Bir(V)=S_3\ltimes(\Z/2)^6,\qquad|\Aut(V)|=384 .
\]
\end{proposition}

\begin{proof}
By adjunction for a smooth complete intersection of four quadrics in $\PP^6$,
$K_V=(K_{\PP^6}+4\cdot\mathcal{O}(2))|_V=(-7+8)\mathcal{O}_V(1)=\mathcal{O}_V(1)$,
which is ample. The Chern computation $c(T_V)\equiv(1+H)^7(1+2H)^{-4}\equiv
1-H+5H^2 \pmod{H^3}$ with $\deg V=H^2\cdot V=2^4=16$ gives $K_V^2=16$,
$c_2(V)=5\cdot16=80$, hence $\chi(\mathcal{O}_V)=(K_V^2+c_2)/12=8$; the Lefschetz
hyperplane theorem gives $q=0$ and so $p_g=\chi-1=7$. Since $K_V$ is ample the
Kodaira dimension is $2$, so $V$ is of general type and minimal, and is not K3
(a K3 surface has $K=0$). The Chern bookkeeping is recorded symbolically in the
project's \texttt{aut\_birV} package.

For the automorphism group: since $K_V=\mathcal{O}_V(1)$ the surface is
canonically embedded, so every automorphism is the restriction of a linear map of
$\PP^6$. A determination of the linear stabilizer of the defining ideal shows its
Lie algebra is one-dimensional (the scalar matrix only), so the linear
automorphism group is $0$-dimensional, hence finite; an explicit enumeration of
the stabilizer gives $S_3\ltimes(\Z/2)^6$ of order $384$ (the six permutations of
$(a,b,c)$, each lifting to a unique permutation of $(d,e,f)$ fixing $g$, and the
seven coordinate sign flips modulo the diagonal scalar, giving $(\Z/2)^6$ of order
$64$). The equality $\Aut(V)=\Bir(V)$ is the statement that for a minimal surface
of general type the natural map $\Aut(V)\to\Bir(V)$ is an isomorphism, a
birational map between minimal models of a non-ruled surface being an
isomorphism. This linear-stabilizer determination is a Gr\"obner-basis ideal
computation carried out in the \texttt{aut\_birV} package; it is symbolic
computer algebra and is \emph{not} reproducible in PARI/GP, so we cite it as an
external computer-algebra input rather than as part of the accompanying PARI
scripts.
\end{proof}

\begin{proposition}\label{prop:descent}
Let $L$ be any ample line bundle on $V$ and $H_L$ an associated Weil height. There
is no birational self-map $\sigma\colon V\dashrightarrow V$ and constant $C>0$ for
which
\[
H_L(\sigma(P))\le H_L(P)-C
\]
holds for all $P$ in a $\sigma$-invariant Zariski-dense subset $S\subseteq V(\Q)$.
In particular, no Markov--Vieta or Fermat-style height-strict infinite descent can
be realized by a birational transformation of $V$.
\end{proposition}

\begin{proof}
By Proposition~\ref{prop:geomV}, $V$ is a smooth projective variety with ample
canonical bundle. By Matsumura's theorem \cite{Matsumura}, such a variety admits
no positive-dimensional connected group of birational transformations; together
with the finiteness of the automorphism group of a variety of general type
(Maehara's extension of Severi's theorem \cite{Maehara,Iitaka}), the group
$\Bir(V)$ is finite, and for a minimal surface of general type it equals
$\Aut(V)$.

Let $\sigma\in\Bir(V)=\Aut(V)$. Since $\Bir(V)$ is finite, $\sigma$ has finite
order $k\ge1$, so $\sigma^{k}=\mathrm{id}$. Fix an ample $L$ and a representative
Weil height $H_L\colon V(\Q)\to\R$. Suppose, for contradiction, that
$H_L(\sigma P)\le H_L(P)-C$ for some $C>0$ and all $P$ in a $\sigma$-invariant
Zariski-dense set $S\subseteq V(\Q)$. Summing this inequality along the finite
orbit $P,\sigma P,\dots,\sigma^{k-1}P$ of any $P\in S$ (each $\sigma^{i}P\in S$ by
$\sigma$-invariance, so the inequality applies at every step) gives
\[
H_L(\sigma^{k}P)\le H_L(P)-kC .
\]
But $\sigma^{k}=\mathrm{id}$, so the left-hand side equals $H_L(P)$, forcing
$kC\le0$---a contradiction, since $k\ge1$ and $C>0$. (A Weil height is well
defined only up to $O(1)$; the argument fixes one representative $H_L$ and
evaluates it at the actual orbit points, where the periodicity $\sigma^{k}P=P$ is
exact, so the $O(1)$ ambiguity never enters.) Hence no birational self-map of $V$
can effect a height-strict infinite descent on $V(\Q)$.
\end{proof}

Proposition~\ref{prop:descent} rests on the folklore engine
``ample $K$ $\Rightarrow$ $\Bir$ finite'' of Matsumura and Maehara
\cite{Matsumura,Maehara,Iitaka}; its content here is the explicit application to
the concrete general-type surface $V$. It explains, intrinsically, why the descent
that succeeds for single equations such as $x^4+y^4=z^2$ (where the ambient is a
curve) or the Markov tree for $x^2+y^2+z^2=3xyz$ (the Vieta involution
$z\mapsto3xy-z$) has no analogue acting on the cuboid surface: the relevant ambient
is a rigid surface of general type with a finite group of birational symmetries.
The $S_3\ltimes(\Z/2)^6$ symmetry merely permutes edge labels and flips signs, so
it reshuffles a candidate cuboid among others with the same space diagonal $g^2$,
never toward a contradiction.

\begin{remark}\label{rem:yelle}
The hypothesis is that the descent be effected \emph{by a birational self-map of}
$V$. This is distinct from the elementary descent of
Yelle~\cite{Yelle}, whose mechanism is a prime-propagation argument on
divisibility data along the space diagonal, not a self-map of $V$;
Proposition~\ref{prop:descent} says nothing about descents of that other kind and
does not bear on Yelle's argument. (Two further internal observations---that a
purely multiplicative reformulation in $\Z[i]$ adds no obstruction beyond the
classical fact that each face diagonal must be composite, and that the
Pila--Zannier unlikely-intersection viewpoint has no special locus to act on for
the fourth, body-diagonal condition---are recorded only to spare future effort and
refute no published claim.)
\end{remark}

The arithmetic side furnishes a second, independent obstruction, this time to an
elementary height inequality. The standard per-fiber elliptic curve attached to a
Pythagorean parameter $q$ is
\[
E_{\mathrm{PCP}}(q):\quad y^2=x(x+1)(x+q^2)=x^3+(1+q^2)x^2+q^2x ,
\]
with full rational $2$-torsion.

\begin{observation}\label{obs:noabs}
For every Pythagorean parameter $q$, the curve $E_{\mathrm{PCP}}(q)$ has purely
multiplicative (Kodaira type $I_n$) reduction at every bad prime, so every
non-archimedean N\'eron local height is nonpositive,
$\lambda_p(P)\in[-\tfrac{N_p}{4}\log p,\,0]$ with $N_p=v_p(\Delta_{\min})$, and
$\sum_p\lambda_p(P)\in[-\tfrac14\log|\Delta_{\min}|,\,0]$. Consequently the
positive part of the canonical height
$\widehat h(P)=\lambda_\infty(P)+\sum_p\lambda_p(P)$ resides entirely in the
archimedean term, and no elementary discriminant-only lower bound
$\widehat h(P)\ge c\log|\Delta_{\min}|$ with $c$ independent of the Szpiro ratio
is available: the multiplicative structure supplies no positive height budget.
\end{observation}

Multiplicativity is the vanishing $v_p(c_4)=0$ at each bad prime, verified on
representative fibers in script~\texttt{09\_local\_heights\_In.gp} (for instance at
$(m,n)=(11,2)$ the bad primes $3,7,11,13,23,73$ each have $v_p(c_4)=0$ and Kodaira
type $I_n$); the local-height range is N\'eron's formula for multiplicative
reduction \cite{SilvermanATAEC}, and the sum bound follows from
$\sum_p N_p\log p=\log|\Delta_{\min}|$. The headline ``no absolute constant'' is, at
bottom, the non-effectivity of Lang's height conjecture
\cite{HindrySilverman,Wagener}, recorded here only to locate the one remaining
elementary handle---a bound on the Szpiro ratio---rather than the height itself.
Observation~\ref{obs:noabs} is consistent with, and complementary to, the rank
data of Section~\ref{sec:curve}: the two rank-zero curves there,
$X_+=\texttt{120a2}$ and $X_-=\texttt{80a1}$, are finite over $\Q$ by Kolyvagin's
theorem \cite{kolyvagin} (analytic rank $0$ with $L(X_+,1)=1.26949\ldots$,
$L(X_-,1)=1.00945\ldots$ nonzero), and it is their $\Q(\sqrt5)$-twists
\texttt{600a2}, \texttt{400a2}, not these curves, that raise the genus-five rank to
its genus (Proposition~\ref{prop:rank}, Remark~\ref{rem:wrongfactors}).

These obstructions complement the recent conditional and elementary work of
Peschmann~\cite{peschmann,peschmann2} and Yelle~\cite{Yelle}: they map a boundary
between methods that transfer to the cuboid locus and methods that, for structural
reasons, do not. None of them resolves the perfect cuboid problem, and we make no
such claim.

\section*{Computational verification}

All identities and curve computations were carried out in PARI/GP
version~2.15.4~\cite{pari}. The accompanying scripts are:
\texttt{01\_genus\_and\_decomposition.gp} (genus of $C$; conductors and
$j$-invariants of the axis quotients; \texttt{ellidentify});
\texttt{02\_genus2\_quotient.gp} and \texttt{03\_Cq\_model\_and\_rank.gp}
(the genus-two quotient, its $a_p$ and full Frobenius characteristic polynomial,
and rank data);
\texttt{04\_genuine\_Cq\_lift.gp} (the rational quotient model and the lift of
its points to the $16$ points of $C$);
\texttt{05\_genuine\_quotient\_rank.gp} and
\texttt{06\_true\_decomposition\_rank.gp} (identification of the twisted factors
\texttt{600a2}, \texttt{400a2}, the five-factor $a_p$ match over $32$ good
primes, and the rank-one computations with explicit generators);
\texttt{07\_known\_points\_and\_cohn.gp} (the $16$ points, the height search,
the Pell orbit $Y_n=L_{2n-1}$, and the square test returning $\{1,4\}$); and
\texttt{08\_caseB\_parametrization.gp} (the face and space-diagonal identities
and the degeneracy bookkeeping); and
\texttt{09\_local\_heights\_In.gp} (Observation~\ref{obs:noabs}: purely
multiplicative reduction and the nonpositive local-height range at each bad prime
on three sample fibers). The Chern-number and automorphism-group inputs of
Proposition~\ref{prop:geomV} are symbolic computer-algebra computations from the
project's \texttt{aut\_birV} package, not reproducible in PARI/GP, and are cited
as external inputs. Each PARI script has a captured \texttt{.out} file.

\section*{Acknowledgements}
The author thanks the maintainers of PARI/GP.

\begin{thebibliography}{99}

\bibitem{cohn}
J.~H.~E. Cohn,
\emph{On square Fibonacci numbers},
J. London Math. Soc. \textbf{39} (1964), 537--540.

\bibitem{cohnlucas}
J.~H.~E. Cohn,
\emph{Square Fibonacci numbers, etc.},
Fibonacci Quart. \textbf{2} (1964), 109--113.

\bibitem{niven}
I.~Niven, H.~S. Zuckerman, and H.~L. Montgomery,
\emph{An Introduction to the Theory of Numbers},
5th ed., John Wiley \& Sons, New York, 1991.

\bibitem{cremona}
J.~E. Cremona,
\emph{Algorithms for modular elliptic curves},
2nd ed., Cambridge University Press, Cambridge, 1997.
Online tables: \url{https://www.lmfdb.org/EllipticCurve/Q/}.

\bibitem{kolyvagin}
V.~A. Kolyvagin,
\emph{Euler systems},
in: The Grothendieck Festschrift, Vol.~II,
Progr. Math. \textbf{87}, Birkh\"auser, Boston, 1990, pp.~435--483.

\bibitem{matson}
R.~D. Matson,
\emph{Results of a computer search for a perfect cuboid},
unpublished manuscript, 2009;
\url{https://www.unsolvedproblems.org/}.

\bibitem{pari}
The PARI~Group,
\emph{PARI/GP version 2.15.4},
Univ. Bordeaux, 2023,
\url{https://pari.math.u-bordeaux.fr/}.

\bibitem{peschmann}
R.~Peschmann,
\emph{Quartic reductions and elliptic obstructions for perfect Euler bricks},
preprint, arXiv:2604.09328 (2026).

\bibitem{peschmann2}
R.~Peschmann,
\emph{A torsion-intersection proof of perfect-cuboid nonexistence on $1072$
explicit master-tuple fibers},
preprint, arXiv:2604.28072 (2026).

\bibitem{yoshida}
T.~Yoshida,
\emph{The relationship between face cuboids and elliptic curves},
preprint, arXiv:2407.09825 (2024).

\bibitem{vanLuijk}
R.~van Luijk,
\emph{On perfect cuboids},
Master's thesis, Universiteit Utrecht, 2000.

\bibitem{Matsumura}
H.~Matsumura,
\emph{On algebraic groups of birational transformations},
Atti Accad. Naz. Lincei Rend. Cl. Sci. Fis. Mat. Natur. (8) \textbf{34} (1963),
151--155.

\bibitem{Maehara}
K.~Maehara,
\emph{A finiteness property of varieties of general type},
Math. Ann. \textbf{262} (1983), 101--123.

\bibitem{Iitaka}
S.~Iitaka,
\emph{Algebraic Geometry: An Introduction to Birational Geometry of Algebraic
Varieties},
Graduate Texts in Mathematics \textbf{76}, Springer, New York, 1982.

\bibitem{SilvermanATAEC}
J.~H. Silverman,
\emph{Advanced Topics in the Arithmetic of Elliptic Curves},
Graduate Texts in Mathematics \textbf{151}, Springer, New York, 1994.

\bibitem{HindrySilverman}
M.~Hindry and J.~H. Silverman,
\emph{Diophantine Geometry: An Introduction},
Graduate Texts in Mathematics \textbf{201}, Springer, New York, 2000.

\bibitem{Wagener}
B.~Wagener,
\emph{About the proof of Lang's height conjecture},
preprint, arXiv:1706.03622 (2017).

\bibitem{Yelle}
S.~Yelle,
\emph{An elementary obstruction to the existence of a perfect cuboid},
preprint, arXiv:2602.00239 (2026).

\end{thebibliography}

\end{document}
